```latex
\documentclass[11pt]{article}

\usepackage{amsmath,amssymb,amsthm}
\usepackage[colorlinks=true,linkcolor=blue,citecolor=blue,urlcolor=blue]{hyperref}

\newtheorem{theorem}{Theorem}[section]
\newtheorem{lemma}[theorem]{Lemma}
\newtheorem{proposition}[theorem]{Proposition}
\newtheorem{corollary}[theorem]{Corollary}
\theoremstyle{definition}
\newtheorem{definition}[theorem]{Definition}
\newtheorem{remark}[theorem]{Remark}

\begin{document}

\title{Prime summands of Dehn surgeries on L-space knots}
\date{}
\maketitle

\begin{abstract}
We prove that any closed $3$--manifold obtained by Dehn surgery on an L-space knot in $S^3$ has at most two prime summands. The key input is the Gordon--Luecke theorem describing reducible Dehn surgeries on knots in $S^3$.
\end{abstract}

\section{Statement of the problem}

\begin{definition}[L-space, L-space knot]
An \emph{L-space} is a rational homology $3$--sphere $Y$ such that
\[
\operatorname{rk}\,\widehat{HF}(Y)=|H_1(Y;\mathbb{Z})|
\]
in the sense of Ozsv\'ath--Szab\'o Heegaard Floer homology \cite{OS04}.
An \emph{L-space knot} is a knot $K\subset S^3$ for which there exists a positive integer $p$ such that the Dehn surgery $S^3_p(K)$ is an L-space.
\end{definition}

\begin{definition}[Dehn surgery notation]
Fix a knot $K\subset S^3$, and let $r\in\mathbb{Q}\cup\{\infty\}$ be a slope on $\partial(S^3\setminus \nu K)$ in the usual meridian--longitude coordinates.
We write $S^3_r(K)$ for the closed $3$--manifold obtained by Dehn filling along slope $r$ (with the convention $S^3_\infty(K)=S^3$).
\end{definition}

\begin{definition}[Reducible, prime, prime summands]
A closed, connected, orientable $3$--manifold $Y$ is \emph{reducible} if it contains an embedded $2$--sphere which does not bound an embedded $3$--ball.
It is \emph{prime} if whenever $Y\cong Y_1\# Y_2$ then $Y_1\cong S^3$ or $Y_2\cong S^3$.
The \emph{prime summands} of $Y$ are the factors appearing in its Kneser--Milnor prime decomposition.
\end{definition}

\begin{theorem}\label{thm:main}
Let $K\subset S^3$ be an L-space knot. Then for every slope $r\in\mathbb{Q}\cup\{\infty\}$, the surgered manifold $S^3_r(K)$ has at most two prime summands.
\end{theorem}

\begin{remark}
The proof does not use the L-space hypothesis: the conclusion holds for \emph{every} knot in $S^3$.
\end{remark}

\section{Prime decomposition preliminaries}

We will use the existence and uniqueness of the prime decomposition.

\begin{theorem}[Kneser--Milnor prime decomposition]\label{thm:kneser-milnor}
Every closed, connected, orientable $3$--manifold $Y$ can be written as a connected sum
\[
Y \cong Y_1 \# \cdots \# Y_n
\]
where each $Y_i$ is prime, and this expression is unique up to order and orientation-preserving homeomorphism of the summands.
\end{theorem}

\begin{proof}
Existence is due to Kneser \cite{Kneser29}. Uniqueness is due to Milnor \cite{Milnor62}.
\end{proof}

We also record the standard implication from irreducibility to primeness.

\begin{lemma}\label{lem:irr-implies-prime}
If $Y$ is a closed, connected, orientable, \emph{irreducible} $3$--manifold, then $Y$ is prime.
\end{lemma}

\begin{proof}
Suppose $Y\cong Y_1\#Y_2$. By definition of connected sum there is an embedded sphere $S\subset Y$ realizing the decomposition. Since $Y$ is irreducible, $S$ bounds an embedded $3$--ball $B\subset Y$. Cutting $Y$ along $S$ yields $Y_1\setminus B^3$ and $Y_2\setminus B^3$, and if $S$ bounds a ball in $Y$ then one of these pieces must be a $3$--ball; hence one of $Y_1,Y_2$ is $S^3$. Thus $Y$ is prime.
\end{proof}

\section{Reducible surgeries on knots in \texorpdfstring{$S^3$}{S3}}

The essential input is the following theorem of Gordon--Luecke.

\begin{theorem}[Gordon--Luecke]\label{thm:GL}
Let $K\subset S^3$ be a nontrivial knot. If $S^3_r(K)$ is reducible, then:
\begin{enumerate}
\item $r\in\mathbb{Z}\setminus\{0\}$ (i.e.\ a nonzero integer slope);
\item $S^3_r(K)$ has \emph{exactly two} prime summands, one of which is a lens space.
\end{enumerate}
\end{theorem}

\begin{proof}[Citation]
This is proved in Gordon--Luecke \cite{GL96}. In particular, \cite{GL96} establishes that a reducible surgery on a nontrivial knot in $S^3$ must be integral and that the resulting reducible manifold is a connected sum of a lens space and an irreducible integer homology sphere; equivalently, it has precisely two prime summands.
\end{proof}

\begin{remark}\label{rem:lens-prime}
Lens spaces are irreducible (hence prime) because each $L(p,q)$ is a quotient of $S^3$ by a free finite cyclic action, and any embedded $2$--sphere lifts to embedded $2$--spheres in $S^3$, each bounding a ball; projecting such a ball shows the original sphere bounds a ball in $L(p,q)$. A standard reference is \cite[Chapter~3]{Hempel}.
\end{remark}

\section{Proof of the main theorem}

\begin{proof}[Proof of Theorem~\ref{thm:main}]
Let $K\subset S^3$ be an L-space knot and let $r\in\mathbb{Q}\cup\{\infty\}$.

If $K$ is the unknot, then $S^3_r(K)$ is either a lens space (when $r\neq 0$) or $S^1\times S^2$ (when $r=0$). In either case the manifold is prime, hence has exactly one prime summand.

Assume henceforth that $K$ is nontrivial. If $S^3_r(K)$ is irreducible, then it is prime by Lemma~\ref{lem:irr-implies-prime}, so it has one prime summand.

If instead $S^3_r(K)$ is reducible, then Theorem~\ref{thm:GL} applies and shows that $S^3_r(K)$ has exactly two prime summands (one a lens space). In particular, it has at most two prime summands.

Thus in all cases $S^3_r(K)$ has at most two prime summands.
\end{proof}

\begin{remark}
As noted above, the argument never uses that $K$ is an L-space knot; the same conclusion holds for any knot in $S^3$.
\end{remark}

\begin{thebibliography}{99}

\bibitem{GL96}
C.~McA.~Gordon and J.~Luecke,
\newblock \emph{Dehn surgeries on knots producing reducible manifolds},
\newblock Topology \textbf{35} (1996), no.~2, 385--409.

\bibitem{Hempel}
J.~Hempel,
\newblock \emph{$3$--Manifolds},
\newblock Annals of Mathematics Studies, No.~86, Princeton University Press, 1976.

\bibitem{Kneser29}
H.~Kneser,
\newblock \emph{Geschlossene Fl{\"a}chen in dreidimensionalen Mannigfaltigkeiten},
\newblock Jahresbericht der Deutschen Mathematiker-Vereinigung \textbf{38} (1929), 248--260.

\bibitem{Milnor62}
J.~Milnor,
\newblock \emph{A unique decomposition theorem for $3$--manifolds},
\newblock Amer.\ J.\ Math.\ \textbf{84} (1962), 1--7.

\bibitem{OS04}
P.~Ozsv\'ath and Z.~Szab\'o,
\newblock \emph{Holomorphic disks and topological invariants for closed three-manifolds},
\newblock Ann.\ of Math.\ (2) \textbf{159} (2004), no.~3, 1027--1158.

\end{thebibliography}

\end{document}
```